\section{Posterior computation}\label{sec:computation}

\subsection{Augmentation}

Write the negative binomial in terms of its log odds. Let $\lambda_{it}=r\,e^{\psi_{it}}$ with
\begin{equation}\label{eq:psi}
\psi_{it}=o_{it}+\alpha+\x_{it}^\top\bbeta+\varepsilon_{it},
\end{equation}
so that the intercept $\alpha$ is on the log-odds scale and the intercept of the mean model is $\alpha+\log r$. The slopes $\bbeta$ and the random components $\varepsilon_{it}$ are unchanged. In this parameterization \eqref{eq:nb} becomes
\begin{equation}\label{eq:nb-logit}
\Pr(y_{it}\mid\psi_{it},r)=\frac{\Gamma(y_{it}+r)}{\Gamma(r)\,y_{it}!}\,\frac{\left(e^{\psi_{it}}\right)^{y_{it}}}{\left(1+e^{\psi_{it}}\right)^{y_{it}+r}} .
\end{equation}
By the P\'olya-Gamma identity of \citet{polson2013bayesian},
\begin{equation}\label{eq:pg}
\frac{\left(e^{\psi}\right)^{y}}{\left(1+e^{\psi}\right)^{y+r}}
=2^{-(y+r)}\,e^{\kappa\psi}\int_0^\infty e^{-\omega\psi^2/2}\,p_{\PG}(\omega\mid y+r,0)\,d\omega,
\qquad \kappa=\frac{y-r}{2},
\end{equation}
where $p_{\PG}(\cdot\mid b,0)$ is the P\'olya-Gamma density. Conditional on $\omega_{it}$, the likelihood in $\psi_{it}$ is proportional to $\exp\{-\omega_{it}(\tilde y_{it}-\psi_{it})^2/2\}$ with working response $\tilde y_{it}=\kappa_{it}/\omega_{it}$. Every coefficient that enters $\psi_{it}$ linearly therefore has a Gaussian conditional likelihood, which is what makes the spike-and-slab updates below available in closed form. For the dispersion we use the Chinese restaurant table (CRT) augmentation \citep{zhou2012lognormal, zhou2015negative}: if $L_{it}\mid y_{it},r\sim\CRT(y_{it},r)$, that is $L_{it}=\sum_{l=1}^{y_{it}}b_l$ with independent $b_l\sim\Bern\{r/(r+l-1)\}$, then a gamma prior on $r$ is conjugate.

\subsection{Why the random draws must be integrated out}\label{sec:collapse}

A direct Gibbs sampler would update $(\delta_j,\sigma_j)$ given the current draws $z_{1j},\dots,z_{nj}$. When $\delta_j=0$ these draws do not enter the likelihood, so their full conditional is the prior and they are pure noise. The evidence for switching the parameter on is then computed with a vector of $n$ random numbers that bears no relation to the data, and with realistic $n$ the probability of a switch is negligible. The sampler would rarely, if ever, discover a random parameter it did not start with. We avoid this by updating the block $(\delta_j,\sigma_j,z_{1j},\dots,z_{nj})$ jointly: first $(\delta_j,\sigma_j)$ from their conditional distribution with the $z_{ij}$ integrated out, then the $z_{ij}$ given $(\delta_j,\sigma_j)$. Because the integral over $z_{ij}$ is Gaussian, the marginal evidence has a closed-form integrand, as the next result shows (derivation in \ref{app:derivations}).

\begin{proposition}[Marginal evidence for a random parameter]\label{prop:marginal}
Fix all quantities except $(\sigma_j,z_{1j},\dots,z_{nj})$ and let $e_{it}=\tilde y_{it}-(\psi_{it}-\sigma_j x_{itj}z_{ij})$ be the working residual that excludes the contribution of the $j$th random component. Define the site-level statistics
\[
q_{ij}=\sum_{t=1}^{T_i}\omega_{it}x_{itj}^2,\qquad b_{ij}=\sum_{t=1}^{T_i}\omega_{it}x_{itj}e_{it}.
\]
Then, with $z_{ij}\sim\Normal(0,1)$ integrated out, the ratio of the conditional likelihood at $\sigma_j=\sigma$ to that at $\sigma_j=0$ is $\exp\{\ell_j(\sigma)\}$, where
\begin{equation}\label{eq:ell}
\ell_j(\sigma)=\sum_{i=1}^{n}\left\{-\frac12\log\left(1+\sigma^2q_{ij}\right)+\frac{\sigma^2b_{ij}^2}{2\left(1+\sigma^2q_{ij}\right)}\right\}.
\end{equation}
Given $\sigma_j=\sigma$, the draws are conditionally independent with
\[
z_{ij}\sim\Normal\left(\frac{\sigma b_{ij}}{1+\sigma^2q_{ij}},\ \frac{1}{1+\sigma^2q_{ij}}\right).
\]
\end{proposition}

The Bayes factor for a random parameter,
\begin{equation}\label{eq:bf-sigma}
\mathrm{BF}^{\sigma}_j=\int_{-\infty}^{\infty}\exp\{\ell_j(\sigma)\}\,\phi(\sigma\mid0,\tau_\sigma^2)\,d\sigma,
\end{equation}
is a one-dimensional integral of a smooth function that is symmetric in $\sigma$, so it is computed accurately on a grid over $\sigma\ge0$. Each evaluation of $\ell_j$ costs $O(n)$ operations once $q_{ij}$ and $b_{ij}$ are formed. The analogous Bayes factor for a fixed effect is available in closed form,
\begin{equation}\label{eq:bf-beta}
\mathrm{BF}^{\beta}_j=\left(1+\tau_\beta^2S_j\right)^{-1/2}\exp\left\{\frac{B_j^2}{2\left(S_j+\tau_\beta^{-2}\right)}\right\},
\end{equation}
where $S_j=\sum_{i,t}\omega_{it}x_{itj}^2$, $B_j=\sum_{i,t}\omega_{it}x_{itj}e^{\beta}_{it}$ and $e^{\beta}_{it}=\tilde y_{it}-(\psi_{it}-x_{itj}\beta_j)$.

\subsection{The sampler}

One sweep of the sampler is summarized in Box~1. Steps 1 and 5 are blocked updates in which some variables are integrated out and then redrawn immediately, which preserves the target distribution \citep{vandyk2008partially}. All draws are exact except that of $|\sigma_j|$, for which we use slice sampling \citep{neal2003slice} on the one-dimensional target in step 5(c), or inverse-CDF sampling on the same grid used for the Bayes factor.

\begin{figure}[!htbp]
\centering
\fbox{\begin{minipage}{0.95\textwidth}
\footnotesize\linespread{1}\selectfont\raggedright
\textbf{Box 1. One sweep of the Gibbs sampler.}
\begin{enumerate}[leftmargin=*,itemsep=2pt]
\item For all $(i,t)$, draw $L_{it}\sim\CRT(y_{it},r)$. Then draw
\[
r\sim\Gammad\Big(a_r+\textstyle\sum_{i,t}L_{it},\ b_r+\sum_{i,t}\log(1+e^{\psi_{it}})\Big),
\]
with the second argument a rate. The P\'olya-Gamma variables are integrated out in this step.
\item For all $(i,t)$, draw $\omega_{it}\sim\PG(y_{it}+r,\psi_{it})$ with the new $r$ and set $\tilde y_{it}=(y_{it}-r)/(2\omega_{it})$.
\item Draw $\alpha$ from its Gaussian full conditional, with precision $\sum_{i,t}\omega_{it}+s_\alpha^{-2}$.
\item For each $j$ in random order, update $(\gamma_j,\beta_j)$:
  \begin{enumerate}[label=(\alph*),leftmargin=*]
  \item compute $S_j$, $B_j$ and $\mathrm{BF}^{\beta}_j$ from \eqref{eq:bf-beta}
  \item set $\gamma_j=1$ with probability $\dfrac{\pi_\beta\,\mathrm{BF}^{\beta}_j\,h_j(1)}{\pi_\beta\,\mathrm{BF}^{\beta}_j\,h_j(1)+(1-\pi_\beta)\,h_j(0)}$, where $h_j(g)=\Pr(\delta_j\mid\gamma_j=g)$ accounts for heredity (under strong heredity, $\delta_j=1$ forces $\gamma_j=1$)
  \item if $\gamma_j=1$ draw $\beta_j\sim\Normal\left(B_j/(S_j+\tau_\beta^{-2}),\ (S_j+\tau_\beta^{-2})^{-1}\right)$, otherwise set $\beta_j=0$.
  \end{enumerate}
\item For each random-eligible $j$ in random order, update $(\delta_j,\sigma_j,z_{\cdot j})$:
  \begin{enumerate}[label=(\alph*),leftmargin=*]
  \item compute $q_{ij}$, $b_{ij}$ and $\mathrm{BF}^{\sigma}_j$ in \eqref{eq:bf-sigma} by one-dimensional quadrature of \eqref{eq:ell}
  \item set $\delta_j=1$ with probability $\dfrac{\pi_j\,\mathrm{BF}^{\sigma}_j}{\pi_j\,\mathrm{BF}^{\sigma}_j+1-\pi_j}$, where $\pi_j=\gamma_j\pi_{\sigma1}+(1-\gamma_j)\pi_{\sigma0}$
  \item if $\delta_j=1$ draw $|\sigma_j|$ from the density proportional to $\exp\{\ell_j(\sigma)\}\phi(\sigma\mid0,\tau_\sigma^2)$ on $\sigma>0$, attach a random sign, and draw each $z_{ij}$ from Proposition~\ref{prop:marginal}. Otherwise set $\sigma_j=0$.
  \end{enumerate}
\item Draw $\pi_\beta$, $\pi_{\sigma1}$ (and $\pi_{\sigma0}$ under weak heredity) from their beta full conditionals. If the slab scales are learned (a sensitivity option), also draw $(\tau_\beta^2,\tau_\sigma^2)$ and their auxiliary variables from inverse-gamma full conditionals.
\end{enumerate}
\end{minipage}}
\end{figure}

\paragraph{Hyperparameter updates} Let $k_\beta=\sum_j\gamma_j$ and let $\mathcal J$ be the set of random-eligible covariates. Then
\begin{align*}
\pi_\beta&\sim\Betad\left(a_\beta+k_\beta,\ b_\beta+p-k_\beta\right),\\
\pi_{\sigma1}&\sim\Betad\Big(a_\sigma+\sum_{j\in\mathcal J}\gamma_j\delta_j,\ b_\sigma+\sum_{j\in\mathcal J}\gamma_j(1-\delta_j)\Big),
\end{align*}
with the analogous update for $\pi_{\sigma0}$ over $\{j\in\mathcal J:\gamma_j=0\}$. In the sensitivity option with half-Cauchy priors $\tau_\beta\sim\mathrm{C}^+(0,A_\beta)$ and $\tau_\sigma\sim\mathrm{C}^+(0,A_\sigma)$, written as scale mixtures of inverse-gamma distributions \citep{makalic2016simple},
\begin{align*}
\tau_\beta^2&\sim\IG\Big(\frac{k_\beta+1}{2},\ \frac{1}{\xi_\beta}+\frac12\sum_{j:\gamma_j=1}\beta_j^2\Big),\\
\xi_\beta&\sim\IG\left(1,\ A_\beta^{-2}+\tau_\beta^{-2}\right),
\end{align*}
and similarly for $\tau_\sigma^2$ with $k_\sigma=\sum_j\delta_j$ and the active $\sigma_j$.

\paragraph{Cost} Steps 1 to 4 cost $O(Np)$ per sweep. Step 5 costs $O(N)$ to form $q_{ij}$ and $b_{ij}$ and $O(Gn)$ for a grid of $G$ points, so a sweep costs $O\{Np+Gn|\mathcal J|\}$. P\'olya-Gamma draws with large $y_{it}+r$ can be replaced by a moment-matched normal approximation if needed. Our implementation is in Python, with the grid evaluation of \eqref{eq:ell} compiled with Numba and P\'olya-Gamma draws from the \texttt{polyagamma} package. We use a coarse grid of 48 points followed by a fine grid of 80 points over the region that holds the mass of the integrand, which reproduces the Bayes factor to five decimals against brute-force integration on 400{,}000 points. Run times are reported in Section~\ref{sec:simulation}.

\paragraph{Convergence} We run several chains from dispersed initial specifications (all excluded, all fixed, all random-eligible covariates random) and monitor the rank-normalized $\widehat R$ of \citet{vehtari2021rank} for $\alpha$, $r$, $\beta_j$, $|\sigma_j|$ and the number of selected fixed and random parameters. Strongly correlated candidates can slow one-at-a-time indicator updates. An optional Metropolis step that swaps the states of two correlated candidates addresses this. \tbd{Decide whether the swap move is needed after the first simulations.}
